\chapter{循环：重复再重复}

\begin{introduction}
  \item 为什么需要循环：从“写一百行”到“写三行”
  \item \texttt{while}、\texttt{do-while}、\texttt{for} 三种循环
  \item C89 的循环变量声明规矩
  \item \texttt{break} 与 \texttt{continue}
  \item 嵌套循环与经典算法：累加、素数、斐波那契、辗转相除
  \item 循环的写法、边界与常见错误
\end{introduction}

\chaptermotto{「聪明人把重复交给机器，把思考留给自己。」}

如果要计算 $1+2+3+\cdots+100$，你可以写一百行 \code{sum = sum + i;}，但没人会这么做。循环就是为“重复做同一件事”而生的。

本章会学三种循环：\code{while}、\code{do-while}、\code{for}。它们能力相同，只是适合的场景不同。学完后，你会用循环解决求和、阶乘、素数、斐波那契、最大公约数等一大批问题。

\section{为什么需要循环}

先看一个笨办法：

\begin{codebox}{sum100-hard.c —— 笨办法：写一百行}
#include <stdio.h>

int main(void)
{
    int sum = 0;

    sum = sum + 1;
    sum = sum + 2;
    sum = sum + 3;
    /* ... 中间省略 96 行 ... */
    sum = sum + 100;

    printf("sum = %d\n", sum);
    return 0;
}
\end{codebox}

这段代码有 100 行几乎一样的语句。如果要求和到 10000，就得写一万行；如果要求到 1 亿，根本没法写。更糟糕的是，只要需求变一下（比如只加偶数），就要改一百处。

循环的思路是：把“重复的动作”写一次，让机器重复执行。

\begin{codebox}{sum100.c —— 用循环求和}
#include <stdio.h>

int main(void)
{
    int i;
    int sum = 0;

    for (i = 1; i <= 100; i++) {
        sum = sum + i;
    }

    printf("1 + 2 + ... + 100 = %d\n", sum);
    return 0;
}
\end{codebox}

\begin{outbox}{运行结果}
1 + 2 + ... + 100 = 5050
\end{outbox}

三行循环，替代了一百行。这就是循环的威力。

\begin{figure}[htbp]
  \centering
  \begin{tikzpicture}[x=1cm,y=1cm,>=Stealth,node distance=7mm]
    \node[terminator] (start) {开始};
    \node[process,below=of start] (init) {i = 1, sum = 0};
    \node[decision,below=9mm of init] (cond) {i $\le$ 100 ?};
    \node[process,below=9mm of cond] (body) {sum = sum + i};
    \node[process,below=7mm of body] (inc) {i = i + 1};
    \node[terminator,right=24mm of body] (end) {结束};

    \draw[flow] (start) -- (init);
    \draw[flow] (init) -- (cond);
    \draw[flow] (cond) -- node[lbl,left]{是} (body);
    \draw[flow] (cond) -- node[lbl,above]{否} (end);
    \draw[flow] (body) -- (inc);
    \draw[flow] (inc.east) -- ++(6mm,0) |- (cond.east);
  \end{tikzpicture}
  \caption{累加循环的流程图：条件成立就重复“加—自增—再判断”}
  \label{fig:sum-loop}
\end{figure}

\section{\texttt{while} 循环}

\subsection{基本语法}

\begin{center}
\fbox{\parbox{0.7\textwidth}{
\ttfamily while (条件) \{ \\
\phantom{xxxx}循环体; \\
\ttfamily \}
}}
\end{center}

执行逻辑：先判断条件，为真就执行循环体，执行完再判断条件……直到条件为假，循环结束。

\begin{codebox}{while.c —— while 倒计时}
#include <stdio.h>

int main(void)
{
    int i = 5;

    while (i > 0) {
        printf("%d\n", i);
        i--;            /* 千万别忘了改变循环变量 */
    }

    printf("发射！\n");
    return 0;
}
\end{codebox}

\begin{outbox}{运行结果}
5
4
3
2
1
发射！
\end{outbox}

\begin{pitfall}[死循环]
如果循环体里没有让条件最终变假的语句，循环会永远执行下去：
\begin{codebox}{infinite.c —— 死循环}
#include <stdio.h>

int main(void)
{
    int i = 1;

    while (i <= 5) {
        printf("%d\n", i);
        /* 忘了 i++，条件永远为真 */
    }
    return 0;
}
\end{codebox}
运行后终端会无限刷 1，只能按 \code{Ctrl+C} 强制结束。
\end{pitfall}

\subsection{用 \texttt{while} 实现累加}

\begin{codebox}{while-sum.c —— while 求和}
#include <stdio.h>

int main(void)
{
    int i = 1;
    int sum = 0;

    while (i <= 100) {
        sum = sum + i;
        i++;
    }

    printf("sum = %d\n", sum);
    return 0;
}
\end{codebox}

\begin{outbox}{运行结果}
sum = 5050
\end{outbox}

写 \code{while} 循环时，按“四步法”检查：

\begin{enumerate}
  \item \textbf{初始化}：循环变量从几开始？（\code{i = 1}）
  \item \textbf{条件}：什么时候继续？（\code{i <= 100}）
  \item \textbf{循环体}：每次重复做什么？（\code{sum += i}）
  \item \textbf{更新}：怎么让循环走向结束？（\code{i++}）
\end{enumerate}

少任何一步，循环都可能不工作或停不下来。

\subsection{循环次数不确定时用 \texttt{while}}

\code{while} 最适合“不知道要循环多少次，只知道什么时候停”的场景。例如不断读入数字，直到读到 0 为止：

\begin{codebox}{read-until-zero.c —— 读到 0 为止}
#include <stdio.h>

int main(void)
{
    int n;
    int sum = 0;

    printf("请输入整数（输入 0 结束）：");

    scanf("%d", &n);
    while (n != 0) {
        sum = sum + n;
        scanf("%d", &n);
    }

    printf("总和 = %d\n", sum);
    return 0;
}
\end{codebox}

\begin{outbox}{运行示例}
请输入整数（输入 0 结束）：3 5 7 0
总和 = 15
\end{outbox}

注意 \code{scanf} 出现了两次：循环前读一次，循环体里再读一次。这是“读取—判断—处理—再读取”的经典模式。

\section{\texttt{for} 循环}

\subsection{基本语法}

\begin{center}
\fbox{\parbox{0.86\textwidth}{
\ttfamily for (初始化; 条件; 更新) \{ \\
\phantom{xxxx}循环体; \\
\ttfamily \}
}}
\end{center}

\code{for} 把 \code{while} 的四步压缩到一行，最适合“次数确定”的循环。

\begin{codebox}{for.c —— for 求和}
#include <stdio.h>

int main(void)
{
    int i;
    int sum = 0;

    for (i = 1; i <= 100; i++) {
        sum = sum + i;
    }

    printf("sum = %d\n", sum);
    return 0;
}
\end{codebox}

\subsection{执行顺序}

\code{for} 的执行顺序是：

\begin{enumerate}
  \item 执行\textbf{初始化}（只执行一次）；
  \item 判断\textbf{条件}，为假则结束；
  \item 执行\textbf{循环体}；
  \item 执行\textbf{更新}；
  \item 回到第 2 步。
\end{enumerate}

\begin{figure}[htbp]
  \centering
  \begin{tikzpicture}[x=1cm,y=1cm,>=Stealth,node distance=6mm]
    \node[process] (init) at (0,3.2) {初始化\\\texttt{i = 1}};
    \node[decision] (cond) at (0,1.6) {条件\\\texttt{i <= 100} ?};
    \node[process] (body) at (0,0) {循环体\\\texttt{sum += i}};
    \node[process] (update) at (0,-1.5) {更新\\\texttt{i++}};
    \node[terminator] (end) at (5,1.6) {结束};

    \draw[flow] (init) -- (cond);
    \draw[flow] (cond) -- node[lbl,left]{真} (body);
    \draw[flow] (cond) -- node[lbl,above]{假} (end);
    \draw[flow] (body) -- (update);
    \draw[flow] (update.west) -- ++(-1.2,0) |- (cond.west);
  \end{tikzpicture}
  \caption{\texttt{for} 的执行顺序：初始化一次，然后“条件—循环体—更新”循环}
  \label{fig:for-order}
\end{figure}

\begin{c89box}[C89：循环变量必须提前声明]
C89 不允许在 \code{for} 的初始化部分声明变量。下面两段代码，左边是 C99 写法，右边才是本书要求的 C89 写法：
\begin{codebox}{wrong-for.c / right-for.c —— 两种写法的对比}
/* C99：不允许（本书不使用） */
for (int i = 0; i < n; i++) {
    ...
}

/* C89：正确写法 */
int i;
for (i = 0; i < n; i++) {
    ...
}
\end{codebox}
把 \code{int i;} 写在函数（或语句块）开头，这是本书所有代码的写法。
\end{c89box}

\subsection{\texttt{for} 的灵活写法}

\code{for} 的三个部分都可以为空或使用逗号表达式：

\begin{codebox}{for-flex.c —— for 的灵活用法}
#include <stdio.h>

int main(void)
{
    int i;
    int j;

    /* 从 10 倒数到 1 */
    for (i = 10; i >= 1; i--) {
        printf("%d ", i);
    }
    printf("\n");

    /* 每次加 2：打印 0 2 4 6 8 */
    for (i = 0; i < 10; i += 2) {
        printf("%d ", i);
    }
    printf("\n");

    /* 逗号表达式：同时控制两个变量 */
    for (i = 0, j = 10; i < j; i++, j--) {
        printf("(%d,%d) ", i, j);
    }
    printf("\n");

    return 0;
}
\end{codebox}

\begin{outbox}{运行结果}
10 9 8 7 6 5 4 3 2 1 
0 2 4 6 8 
(0,10) (1,9) (2,8) (3,7) (4,6) 
\end{outbox}

\begin{pitfall}[分号陷阱]
\code{for (i = 0; i < 10; i++);} 末尾的分号会让循环体变成\textbf{空语句}，循环空转十次，后面的代码只执行一次。这是很常见的错误。
\end{pitfall}

\section{\texttt{do-while} 循环}

\code{do-while} 的特点是：\textbf{先执行一次循环体，再判断条件}。所以循环体至少执行一次。

\begin{center}
\fbox{\parbox{0.78\textwidth}{
\ttfamily do \{ \\
\phantom{xxxx}循环体; \\
\ttfamily \} while (条件);
}}
\end{center}

\begin{codebox}{do-while.c —— 至少执行一次}
#include <stdio.h>

int main(void)
{
    int n;

    do {
        printf("请输入 1-10 之间的数：");
        scanf("%d", &n);
    } while (n < 1 || n > 10);

    printf("你输入了 %d\n", n);
    return 0;
}
\end{codebox}

\begin{outbox}{运行示例}
请输入 1-10 之间的数：0
请输入 1-10 之间的数：11
请输入 1-10 之间的数：5
你输入了 5
\end{outbox}

这种“至少问一次，不合法就再问”的场景，\code{do-while} 比 \code{while} 更自然。

\begin{codebox}{digits.c —— 数一数有几位}
#include <stdio.h>

int main(void)
{
    int n;
    int count = 0;

    printf("请输入一个非负整数：");
    scanf("%d", &n);

    do {
        count++;
        n = n / 10;
    } while (n != 0);

    printf("共有 %d 位\n", count);
    return 0;
}
\end{codebox}

\begin{outbox}{运行示例}
请输入一个非负整数：2026
共有 4 位
\end{outbox}

如果 \code{n} 是 0，\code{do-while} 也能正确输出 1 位；换成 \code{while (n != 0)} 则一次都不会执行，结果为 0。

\section{\texttt{break} 与 \texttt{continue}}

\subsection{\texttt{break}：提前跳出循环}

\code{break} 立即结束当前循环，跳到循环后面的语句。

\begin{codebox}{break.c —— 找到就停}
#include <stdio.h>

int main(void)
{
    int i;

    for (i = 1; i <= 100; i++) {
        if (i * i > 50) {
            printf("第一个平方大于 50 的数是 %d\n", i);
            break;
        }
    }

    return 0;
}
\end{codebox}

\begin{outbox}{运行结果}
第一个平方大于 50 的数是 8
\end{outbox}

\code{break} 也常用于“搜索”：找到目标就停止，不必把剩下的都检查一遍。

\subsection{\texttt{continue}：跳过本次，继续下一次}

\code{continue} 跳过循环体剩下的部分，直接进入下一次循环（对 \code{for} 来说，会先执行更新部分）。

\begin{codebox}{continue.c —— 只打印奇数}
#include <stdio.h>

int main(void)
{
    int i;

    for (i = 1; i <= 10; i++) {
        if (i % 2 == 0) {
            continue;       /* 偶数跳过 */
        }
        printf("%d ", i);
    }
    printf("\n");

    return 0;
}
\end{codebox}

\begin{outbox}{运行结果}
1 3 5 7 9 
\end{outbox}

\begin{pitfall}[\texttt{break} 只跳出最内层循环]
在嵌套循环里，\code{break} 只结束它所在的那一层，外层循环继续。如果要一次跳出多层，需要额外的标志变量或 \code{goto}（不推荐）。
\end{pitfall}

\begin{table}[htbp]
  \centering
  \caption{\texttt{break} 与 \texttt{continue} 的区别}
  \begin{tabular}{@{}lll@{}}
    \toprule
    语句 & 作用 & 之后执行什么 \\
    \midrule
    \code{break} & 立即结束当前循环 & 循环后面的语句 \\
    \code{continue} & 跳过本次剩余部分 & 下一次循环（\code{for} 先更新） \\
    \bottomrule
  \end{tabular}
\end{table}

\section{嵌套循环}

循环体里再写一个循环，叫\textbf{嵌套循环}。外层循环每执行一次，内层循环完整地执行一轮。

\begin{codebox}{nested-loop.c —— 嵌套循环}
#include <stdio.h>

int main(void)
{
    int i, j;

    for (i = 1; i <= 3; i++) {
        for (j = 1; j <= 4; j++) {
            printf("(%d,%d) ", i, j);
        }
        printf("\n");
    }

    return 0;
}
\end{codebox}

\begin{outbox}{运行结果}
(1,1) (1,2) (1,3) (1,4) 
(2,1) (2,2) (2,3) (2,4) 
(3,1) (3,2) (3,3) (3,4) 
\end{outbox}

外层 \code{i} 每取一个值，内层 \code{j} 就从 1 到 4 走一遍，所以共执行 $3 \times 4 = 12$ 次。

\subsection{例：九九乘法表}

\begin{example}
输出九九乘法表（下三角）。
\end{example}
\begin{solution}
思路如下：

\begin{codebox}{multiply.c —— 九九乘法表}
#include <stdio.h>

int main(void)
{
    int i, j;

    for (i = 1; i <= 9; i++) {
        for (j = 1; j <= i; j++) {
            printf("%d*%d=%-2d ", j, i, i * j);
        }
        printf("\n");
    }

    return 0;
}
\end{codebox}

\begin{outbox}{运行结果（节选）}
1*1=1  
1*2=2  2*2=4  
1*3=3  2*3=6  3*3=9  
...
1*9=9  2*9=18 3*9=27 ... 9*9=81
\end{outbox}

内层循环的上限是 \code{i}，所以每行输出 \code{i} 项，形成下三角。\texttt{\%-2d} 表示左对齐、占两格，让各列大致对齐。
\end{solution}

\subsection{例：星号三角形}

\begin{example}
输入行数 $n$，输出一个由星号组成的直角三角形。
\end{example}
\begin{solution}
思路如下：

\begin{codebox}{stars.c —— 星号三角形}
#include <stdio.h>

int main(void)
{
    int n;
    int i, j;

    printf("请输入行数：");
    scanf("%d", &n);

    for (i = 1; i <= n; i++) {
        for (j = 1; j <= i; j++) {
            printf("*");
        }
        printf("\n");
    }

    return 0;
}
\end{codebox}

\begin{outbox}{运行示例（n = 5）}
*
**
***
****
*****
\end{outbox}

如果要做“金字塔”（居中对齐），需要在星号前打印空格：第 \code{i} 行先打印 \code{n - i} 个空格，再打印 \code{2*i - 1} 个星号。这个练习留给你自己完成。
\end{solution}

\section{循环的经典算法}

\subsection{累加与累乘}

累加的模式是 \code{sum = sum + 项}，累乘的模式是 \code{product = product * 项}。注意累乘的初始值必须是 1，累加的初始值必须是 0。

\begin{codebox}{factorial.c —— 阶乘}
#include <stdio.h>

int main(void)
{
    int n;
    int i;
    int fact = 1;

    printf("请输入 n：");
    scanf("%d", &n);

    for (i = 1; i <= n; i++) {
        fact = fact * i;
    }

    printf("%d! = %d\n", n, fact);
    return 0;
}
\end{codebox}

\begin{outbox}{运行示例}
请输入 n：10
10! = 3628800
\end{outbox}

\begin{pitfall}[阶乘增长极快]
$13! = 6227020800$，已经超过 32 位 \code{int} 的范围。计算更大的阶乘需要 \code{long long}（C99）或自行实现大整数。不要假设 \code{int} 能装下所有结果。
\end{pitfall}

\subsection{最大值与最小值}

\begin{codebox}{max-min.c —— 求最大值和最小值}
#include <stdio.h>

int main(void)
{
    int n, i;
    int x;
    int max, min;

    printf("请输入数据个数：");
    scanf("%d", &n);

    printf("请输入 %d 个整数：", n);
    scanf("%d", &x);
    max = x;
    min = x;

    for (i = 2; i <= n; i++) {
        scanf("%d", &x);
        if (x > max) {
            max = x;
        }
        if (x < min) {
            min = x;
        }
    }

    printf("最大值 = %d，最小值 = %d\n", max, min);
    return 0;
}
\end{codebox}

\begin{outbox}{运行示例}
请输入数据个数：5
请输入 5 个整数：3 9 1 7 5
最大值 = 9，最小值 = 1
\end{outbox}

注意第一个数要先读进来作为 \code{max} 和 \code{min} 的初始值，而不是把初值设成 0——否则全是负数时最大值会错。

\subsection{判断素数}

素数是只能被 1 和自身整除的大于 1 的整数。最直接的判断方法是试除：

\begin{codebox}{prime.c —— 判断素数}
#include <stdio.h>

int main(void)
{
    int n;
    int i;
    int is_prime = 1;

    printf("请输入一个整数：");
    scanf("%d", &n);

    if (n <= 1) {
        is_prime = 0;
    } else {
        for (i = 2; i * i <= n; i++) {
            if (n % i == 0) {
                is_prime = 0;
                break;
            }
        }
    }

    if (is_prime) {
        printf("%d 是素数\n", n);
    } else {
        printf("%d 不是素数\n", n);
    }

    return 0;
}
\end{codebox}

\begin{outbox}{运行示例}
请输入一个整数：97
97 是素数
\end{outbox}

为什么循环条件是 \code{i * i <= n}？因为如果 $n$ 有大于 $\sqrt{n}$ 的因子，那么它一定还有一个小于 $\sqrt{n}$ 的因子。检查到 $\sqrt{n}$ 就够了。

\subsection{打印 100 以内的素数}

\begin{codebox}{primes100.c —— 100 以内的素数}
#include <stdio.h>

int main(void)
{
    int n, i, is_prime;

    for (n = 2; n <= 100; n++) {
        is_prime = 1;
        for (i = 2; i * i <= n; i++) {
            if (n % i == 0) {
                is_prime = 0;
                break;
            }
        }
        if (is_prime) {
            printf("%d ", n);
        }
    }
    printf("\n");

    return 0;
}
\end{codebox}

\begin{outbox}{运行结果}
2 3 5 7 11 13 17 19 23 29 31 37 41 43 47 53 59 61 67 71 73 79 83 89 97
\end{outbox}

这是嵌套循环的典型用法：外层枚举候选数，内层判断素数。

\subsection{斐波那契数列}

斐波那契数列：$F_1 = 1$，$F_2 = 1$，$F_n = F_{n-1} + F_{n-2}$。

\begin{codebox}{fib.c —— 前 20 项斐波那契数}
#include <stdio.h>

int main(void)
{
    int i;
    int a = 1;
    int b = 1;
    int c;

    printf("%d %d ", a, b);
    for (i = 3; i <= 20; i++) {
        c = a + b;
        printf("%d ", c);
        a = b;
        b = c;
    }
    printf("\n");

    return 0;
}
\end{codebox}

\begin{outbox}{运行结果}
1 1 2 3 5 8 13 21 34 55 89 144 233 377 610 987 1597 2584 4181 6765
\end{outbox}

这个“滚动变量”的技巧很常用：只保留最近两个值，不断往前推。

\subsection{辗转相除法求最大公约数}

\begin{codebox}{gcd.c —— 最大公约数}
#include <stdio.h>

int main(void)
{
    int a, b;
    int t;

    printf("请输入两个正整数：");
    scanf("%d%d", &a, &b);

    while (b != 0) {
        t = a % b;
        a = b;
        b = t;
    }

    printf("最大公约数 = %d\n", a);
    return 0;
}
\end{codebox}

\begin{outbox}{运行示例}
请输入两个正整数：48 18
最大公约数 = 6
\end{outbox}

辗转相除法（欧几里得算法）的思想是：$\gcd(a, b) = \gcd(b, a \bmod b)$，直到余数为 0，此时的除数就是最大公约数。

\subsection{水仙花数}

水仙花数是指一个三位数，其各位数字的立方和等于它本身，例如 $153 = 1^3 + 5^3 + 3^3$。

\begin{codebox}{narcissus.c —— 水仙花数}
#include <stdio.h>

int main(void)
{
    int n;
    int a, b, c;

    for (n = 100; n <= 999; n++) {
        a = n / 100;
        b = n / 10 % 10;
        c = n % 10;
        if (a * a * a + b * b * b + c * c * c == n) {
            printf("%d ", n);
        }
    }
    printf("\n");

    return 0;
}
\end{codebox}

\begin{outbox}{运行结果}
153 370 371 407
\end{outbox}

\subsection{反转整数}

\begin{codebox}{reverse-int.c —— 反转整数}
#include <stdio.h>

int main(void)
{
    int n;
    int rev = 0;

    printf("请输入一个整数：");
    scanf("%d", &n);

    while (n != 0) {
        rev = rev * 10 + n % 10;
        n = n / 10;
    }

    printf("反转后 = %d\n", rev);
    return 0;
}
\end{codebox}

\begin{outbox}{运行示例}
请输入一个整数：12345
反转后 = 54321
\end{outbox}

这个算法对负数也能工作，但反转后的值可能溢出（例如反转 2147483647），实际工程中要检查。

\section{怎样写好一个循环}

\subsection{四步检查法}

写循环前，先回答四个问题：

\begin{enumerate}
  \item \textbf{从哪里开始}？（初始化）
  \item \textbf{做到什么时候停}？（条件）
  \item \textbf{每次做什么}？（循环体）
  \item \textbf{怎么走向结束}？（更新）
\end{enumerate}

写完后，用三个值测试：\textbf{最小规模}（0 次或 1 次）、\textbf{典型值}、\textbf{边界值}。例如求和 1 到 100，要测试 \code{n = 0}、\code{n = 1}、\code{n = 100}。

\subsection{边界：\texttt{<=} 还是 \texttt{<}}

\begin{center}
\begin{tabular}{@{}lll@{}}
\toprule
写法 & 循环次数 & 说明 \\
\midrule
\code{for (i = 1; i <= n; i++)} & $n$ 次 & 包含 $n$ \\
\code{for (i = 0; i < n; i++)} & $n$ 次 & 不包含 $n$，适合数组下标 \\
\code{for (i = 1; i < n; i++)} & $n-1$ 次 & 少一次，常犯的 off-by-one 错误 \\
\bottomrule
\end{tabular}
\end{center}

数组下标从 0 开始时，\code{for (i = 0; i < n; i++)} 是最标准的写法（第 6 章会反复用）。

\subsection{循环不变量}

\textbf{循环不变量}是每次循环开始时都成立的性质。例如求和循环的不变量是：

\begin{center}
\kaishu 在每轮循环开始时，\code{sum} 等于 $1 + 2 + \cdots + (i-1)$。
\end{center}

验证一下：初始 \code{i = 1} 时，\code{sum = 0}，空和确实为 0；每轮把 \code{i} 加进 \code{sum}，然后 \code{i} 加 1，不变量继续成立；循环结束时 \code{i = 101}，\code{sum} 就是 $1 + \cdots + 100$。

养成想“不变量”的习惯，能帮你快速检查循环写得对不对。

\section{循环的常见错误}

\begin{center}
\begin{tabular}{@{}p{0.38\textwidth}p{0.52\textwidth}@{}}
\toprule
错误 & 后果 \\
\midrule
忘了更新循环变量 & 死循环 \\
\code{for (...);} 多写分号 & 循环体为空，后面的语句只执行一次 \\
\code{while (i <= n)} 里 \code{n} 被修改 & 循环次数不可预测 \\
用浮点数做循环变量 & 累积误差可能让循环多一次或少一次 \\
\code{unsigned} 变量写 \code{i >= 0} & 永远为真，死循环 \\
循环体内修改循环变量 & 次数混乱 \\
嵌套循环里 \code{break} & 只跳出内层 \\
\bottomrule
\end{tabular}
\end{center}

\begin{pitfall}[无符号变量的倒序循环]
\begin{codebox}{unsigned-loop.c —— 死循环陷阱}
#include <stdio.h>

int main(void)
{
    unsigned int i;

    for (i = 10; i >= 0; i--) {
        printf("%u ", i);   /* 0 之后再减 1 会变成 4294967295 */
    }
    return 0;
}
\end{codebox}
这个循环停不下来。正确写法是用 \code{int}，或者写成 \code{for (i = 10; i > 0; i--)} 并单独处理 0。
\end{pitfall}

\begin{pitfall}[浮点数循环变量]
\begin{codebox}{float-loop.c —— 浮点循环的隐患}
for (x = 0.0; x != 1.0; x += 0.1) {
    ...
}
\end{codebox}
0.1 在二进制里不能精确表示，累加十次后 \code{x} 可能不等于 1.0，循环停不下来。循环次数应该用整数控制。
\end{pitfall}

\section{综合例题}

\begin{example}
输入 $n$，计算 $1 + \frac{1}{2} + \frac{1}{3} + \cdots + \frac{1}{n}$，保留 4 位小数。
\end{example}
\begin{solution}
思路如下：

\begin{codebox}{harmonic.c —— 调和级数}
#include <stdio.h>

int main(void)
{
    int n, i;
    double sum = 0.0;

    printf("请输入 n：");
    scanf("%d", &n);

    for (i = 1; i <= n; i++) {
        sum = sum + 1.0 / i;
    }

    printf("sum = %.4f\n", sum);
    return 0;
}
\end{codebox}

注意 \code{1.0 / i} 而不是 \code{1 / i}。后者是整数除法，除了 \code{i = 1} 以外全是 0。
\end{solution}

\begin{example}
输入若干个整数（以 0 结束），统计正数、负数的个数，并计算平均值。
\end{example}
\begin{solution}
思路如下：

\begin{codebox}{stat.c —— 统计正负数}
#include <stdio.h>

int main(void)
{
    int n;
    int positive = 0;
    int negative = 0;
    int sum = 0;
    int count = 0;

    printf("请输入整数（0 结束）：");
    scanf("%d", &n);

    while (n != 0) {
        if (n > 0) {
            positive++;
        } else {
            negative++;
        }
        sum = sum + n;
        count++;
        scanf("%d", &n);
    }

    printf("正数 %d 个，负数 %d 个\n", positive, negative);
    if (count > 0) {
        printf("平均值 = %.2f\n", (double)sum / count);
    } else {
        printf("没有输入有效数据\n");
    }

    return 0;
}
\end{codebox}

这里用 \code{(double)sum / count} 做浮点除法；同时处理了“一个数都没输入”的情况，避免除以 0。
\end{solution}

\begin{example}
输出 1000 以内的所有完数。完数是指一个数恰好等于它的真因子之和，例如 $6 = 1 + 2 + 3$。
\end{example}
\begin{solution}
思路如下：

\begin{codebox}{perfect.c —— 完数}
#include <stdio.h>

int main(void)
{
    int n, i, sum;

    for (n = 2; n <= 1000; n++) {
        sum = 0;
        for (i = 1; i < n; i++) {
            if (n % i == 0) {
                sum = sum + i;
            }
        }
        if (sum == n) {
            printf("%d ", n);
        }
    }
    printf("\n");

    return 0;
}
\end{codebox}

\begin{outbox}{运行结果}
6 28 496
\end{outbox}
\end{solution}

\begin{example}
输入一个整数，输出它的二进制表示（从最高位到最低位）。
\end{example}
\begin{solution}
可以用“掩码”从高位往低位检查。以 32 位 \code{unsigned int} 为例：

\begin{codebox}{to-binary.c —— 输出二进制}
#include <stdio.h>

int main(void)
{
    unsigned int n;
    int i;

    printf("请输入一个非负整数：");
    scanf("%u", &n);

    for (i = 31; i >= 0; i--) {
        if (n & (1u << i)) {
            printf("1");
        } else {
            printf("0");
        }
        if (i % 8 == 0) {
            printf(" ");
        }
    }
    printf("\n");

    return 0;
}
\end{codebox}

\begin{outbox}{运行示例}
请输入一个非负整数：182
00000000 00000000 00000000 10110110 
\end{outbox}

\code{1u << i} 生成只有第 \code{i} 位为 1 的掩码，\texttt{n \& 掩码} 判断该位是否为 1。这是位运算和循环的经典结合。
\end{solution}

\section{拔高：循环的效率与递归}

\subsection{时间复杂度直觉}

循环执行次数决定了程序要花多少时间。求 1 到 $n$ 的和需要循环 $n$ 次，记作 $O(n)$；判断一个数是否为素数只需要试除到 $\sqrt{n}$，记作 $O(\sqrt{n})$。数据量增大时，$O(\sqrt{n})$ 比 $O(n)$ 快得多。

\begin{center}
\begin{tabular}{@{}lll@{}}
\toprule
$n$ & $O(n)$ 次操作 & $O(\sqrt{n})$ 次操作 \\
\midrule
100 & 100 & 10 \\
10000 & 10000 & 100 \\
1000000 & 1000000 & 1000 \\
\bottomrule
\end{tabular}
\end{center}

\subsection{递归：循环的另一种写法}

有些问题用“自己调用自己”描述更自然，例如阶乘：

\begin{codebox}{fact-rec.c —— 递归求阶乘}
#include <stdio.h>

int factorial(int n)
{
    if (n <= 1) {
        return 1;
    }
    return n * factorial(n - 1);
}

int main(void)
{
    printf("%d\n", factorial(5));
    return 0;
}
\end{codebox}

递归和循环可以互相转换。初学阶段先把循环写熟，递归等学完函数（后续章节）再深入。

\subsection{\texttt{goto}：不推荐，但要知道}

\code{goto} 可以跳到同一函数内的标号处，常用于“多层循环一次跳出”或“统一清理资源”。但它会破坏代码结构，现代工程中很少使用。知道有这个东西即可，不要主动用。

\section{本章小结}

\begin{center}
\begin{tabular}{@{}p{0.24\textwidth}p{0.66\textwidth}@{}}
\toprule
主题 & 要点 \\
\midrule
\code{while} & 先判断后执行；适合次数不确定的循环 \\
\code{for} & 初始化、条件、更新集中在一行；适合次数确定的循环 \\
\code{do-while} & 先执行一次再判断；适合“至少做一次”的场景 \\
C89 循环变量 & 必须在循环外声明 \\
\code{break} & 立即结束当前循环 \\
\code{continue} & 跳过本次，进入下一次 \\
嵌套循环 & 外层每走一步，内层走完整一轮 \\
写循环 & 初始化、条件、循环体、更新，四步缺一不可 \\
常见错误 & 死循环、多余分号、边界差一、无符号倒序 \\
\bottomrule
\end{tabular}
\end{center}

\begin{tipbox}[下一章预告]
有了循环，我们可以处理大量数据了。但数据放哪里？如果每读一个数就丢了，循环再厉害也没用。下一章学习\textbf{数组}——把一堆同类型的数据整齐地存在一起。
\end{tipbox}

\begin{problemset}
  \item 用 \code{for}、\code{while}、\code{do-while} 三种循环分别求 $1+2+\cdots+100$。
  \item 输入 $n$，求 $n!$，并测试 $n=10$、$n=13$ 的结果，解释差异。
  \item 输入 10 个整数，求最大值、最小值和平均值。
  \item 输出 1\textasciitilde100 中所有能被 3 或 5 整除的数，每行 5 个。
  \item 判断一个数是否为素数，输出 100 以内的所有素数。
  \item 输出前 30 项斐波那契数列。
  \item 输入两个正整数，用辗转相除法求最大公约数和最小公倍数。
  \item 输出所有三位水仙花数。
  \item 输入一个整数，输出它的位数、各位数字之和、逆序数。
  \item 思考题：为什么 \code{for (i = 0; i <= n; i++)} 在处理长度为 \code{n} 的数组时会越界？应该怎么改？
\end{problemset}
